% ===== 原文 pp.18–19 =====
% tr-02c：定理 6、定理 7、定理 8（p.18，编号顺序 6→7→8 不可错位）与长度恒等式的证明开头（p.19）
% 编号公式 (1)(2)(3) 全在 p.18：(1)=H 的定义；(2)=D、R 的定义（两行共用编号）；
% (3)=∂D/∂x 与 2∂R/∂x 两条关系（两行共用编号）。
% 按 Lead 裁决：用 plain TeX display + \eqno 复刻原文右对齐编号（禁止 \tag、禁止带编号环境）。

% [原文 p.18]
\begin{thm}{6}
（Tan--Wong--Zhang [TWZ06]，以及 Mirzakhani [Mir07a]，二者都沿袭 McShane [McS98]。）边界点有一个 Cantor 集划分，由正交出发的测地线的行为给出：
$$
\begin{gathered}
\beta=\{\text{与边界正交且只相交一次的简单测地线}\}\\
\cup\{\text{正交边界测地线}\}\cup\{\text{边界斜交的简单测地线}\}\cup\{\text{非简单测地线}\}.
\end{gathered}
$$
\end{thm}

以 $\epsilon$ 为半径的、简单完全测地线（在各交点处与边界正交）的邻域，是可数多条细走廊（corridors）的并。在万有覆盖中，这些走廊由基本群中的约化双无穷词（reduced bi infinite words）描述。通过分析走廊的数目与宽度，Birman--Series 表明该集合非常稀薄。

\begin{thm}{7}
（Birman--Series，[BS85]。）简单测地线的测度为零。简单完全测地线（在各交点处与边界正交）之集 $S$ 的 Hausdorff 维数为 $1$。$S$ 与边界的交集，其 Hausdorff 维数与测度均为 $0$。
\end{thm}

长度恒等式的基本求和项是一个有理指数函数。在 $\R^2$ 上定义函数 $H$ 为
$$
H(x,y)=\frac{1}{1+e^{\frac{x+y}{2}}}+\frac{1}{1+e^{\frac{x-y}{2}}}\eqno{(1)}
$$
并在 $\R^3$ 上定义相应的函数 $\mathcal{D},\mathcal{R}$ 为
$$
\begin{aligned}
\mathcal{D}(x,y,z)&=2\log\left(\frac{e^{\frac{x}{2}}+e^{\frac{y+z}{2}}}{e^{\frac{-x}{2}}+e^{\frac{y+z}{2}}}\right)\quad\text{以及}\\
\mathcal{R}(x,y,z)&=x-\log\left(\frac{\cosh\frac{y}{2}+\cosh\frac{x+z}{2}}{\cosh\frac{y}{2}+\cosh\frac{x-z}{2}}\right).
\end{aligned}
\eqno{(2)}
$$

% [Tutor Note: (2) 的 \(\mathcal{D}\) 分子分母第二项都是 \(e^{(y+z)/2}\)（不是 \(e^{(x+z)/2}\)），这是最易抄错处：
%  正因为分母里出现的是 \(e^{-x/2}\) 与同一个 \(e^{(y+z)/2}\)，\(\mathcal{D}\) 才可能写成
%  \(\mathcal D(x,y,z)=2\log\big(\frac{e^{x/2}+A}{e^{-x/2}+A}\big)\) 且 \(A=e^{(y+z)/2}\) 与 \(x\) 无关；
%  核对量纲：把 \(x=0\) 代入得 \(\mathcal D(0,y,z)=2\log\big(\frac{1+A}{1+A}\big)=0\) 对一切 \(y,z\) 成立。
%  \(\mathcal{R}\) 里两个 \(\cosh(y/2)\) 相同、只把 \(\cosh((x+z)/2)\) 换成 \(\cosh((x-z)/2)\)，
%  与 p.19 的 \(2\ell_{\beta_\alpha}=\ell_\beta-\mathcal R(\ell_\beta,\ell_\lambda,\ell_\alpha)\) 的第 2、3 变量顺序自洽。]

函数 $\mathcal{D}$ 与 $\mathcal{R}$ 与 $H$ 的关系如下，
$$
\begin{aligned}
\frac{\partial}{\partial x}\mathcal{D}(x,y,z)&=H(y+z,x),\qquad\mathcal{D}(0,0,0)=0\quad\text{以及}\\
2\frac{\partial}{\partial x}\mathcal{R}(x,y,z)&=H(z,x+y)+H(z,x-y),\qquad\mathcal{R}(0,0,0)=0.
\end{aligned}
\eqno{(3)}
$$

% [Tutor Note: (3) 可以当作 \(\mathcal D,\mathcal R\) 的积分定义，用于抄写自检，三个可验算的锚点：
%  ① 由 \(\partial_x\mathcal D=H(y+z,x)\) 与 \(\mathcal D(0,0,0)=0\) 得 \(\mathcal D(x,y,z)=\int_0^xH(y+z,t)\dd t\)，
%     故 \(\mathcal D(x,0,0)=2\log e^{x/2}=x\)——这正是 McShane 恒等式收敛的来源：短测地线（\(\ell_\alpha\to0\)）
%     在 \(\sum_\alpha\mathcal D(L,\ell_\alpha,\ell_\alpha)\) 中的贡献趋于 \(L\) 本身（p.19 的亏格 1 情形）。
%  ② \(H(0,y)\equiv1\)：因为 \(\frac{1}{1+e^{y/2}}+\frac{1}{1+e^{-y/2}}=\frac{1}{1+e^{y/2}}+\frac{e^{y/2}}{1+e^{y/2}}=1\)。
%  ③ 于是 (3) 第二式在 \(y=z=0\) 时化为 \(2\partial_x\mathcal R(x,0,0)=H(0,x)+H(0,x)=2\)，
%     配合 \(\mathcal R(0,0,0)=0\) 得 \(\mathcal R(x,0,0)=x\)，与直接代入 \(\mathcal R\) 的定义
%     （此时分式的分子分母同为 \(\cosh0+\cosh\frac{x}{2}\)，对数为 0）完全一致。若把 (2) 中
%     \(\mathcal R\) 分母的 \(\cosh\frac{x-z}{2}\) 误抄成 \(\cosh\frac{x+z}{2}\)，则 \(\mathcal R\equiv x\)，(3) 立刻不成立。]

\begin{thm}{8}
[Mir07a, 定理 1.3 与 4.2] 以及 [TWZ06, 定理 1.8]。Mirzakhani--McShane 恒等式。设 $R$ 是双曲曲面，其边界为 $\beta_j$、长度为 $L_j$，则
$$
L_1=\sum_{\alpha_1,\alpha_2}\mathcal{D}(L_1,\ell_{\alpha_1}(R),\ell_{\alpha_2}(R))+\sum_{j=2}^n\sum_{\alpha}\mathcal{R}(L_1,L_j,\ell_\alpha(R)),
$$
其中第一个和遍历所有无序的简单闭测地线对，使 $\beta_1,\alpha_1,\alpha_2$ 围出一条嵌入的裤子；二重和遍历简单闭测地线，使 $\beta_1,\beta_j,\alpha$ 围出一条嵌入的裤子。
\end{thm}

% [Tutor Note: 定理 8 的两个和都遍历"嵌入的裤子"而非所有简单闭测地线，这是长度恒等式与一般求和最不同之处：
%  对给定的一条 \(\beta\)-袖口裤子，它在 \(\beta_1\) 上恰好占一段主间隙（p.19 的证明开头即用此），
%  所以两个和加起来正是 \(\beta_1\) 上 Cantor 集补区长度的总和，等于 \(L_1\)。
%  注意二重和的下标从 \(j=2\) 起：左端已固定 \(\beta_1\)，第二个和只把它与其余边界 \(\beta_j\)（\(j\ge2\)）配对；\(j=1\) 等于把 \(\beta_1\) 与自身凑成一条裤子的两条边界，不是嵌入的裤子。
%  另外计数上要小心"无序对"\(\{\alpha_1,\alpha_2\}\)：若 \(\alpha_1=\alpha_2\)（即该裤子的两条 \(\alpha\) 边界相同，
%  对应 \((g,n)=(1,1)\) 这类特殊情形），一对只算一次，这与 p.19 用 \(\mathcal D(L,\ell_\alpha,\ell_\alpha)\) 写环面情形一致。]

% [原文 p.19]
\begin{proof}
由上面各定理，$\ell_\beta$ 等于遍历嵌入的 $\beta$-袖口裤子的主间隙长度之和，再加上双边界袖口裤子（double boundary cuff pants）的相应长度。为求出主间隙长度，先给出补区间长度的一个公式。在一条裤子中，从 $\beta$ 到 $\alpha$ 的正交边界测地线 $\delta_\alpha$ 平分一个补区间（见图 6）。设 $\beta_\alpha$ 是所得两个半区间之一。线段 $\beta_\alpha$ 的一端有测地线 $\delta_\beta$ 出发，另一端有一条通向 $\alpha$ 的螺线 $\sigma$ 出发。在万有覆盖中，考虑相邻的提升 $\wt{\delta}_\alpha$、$\wt{\beta}_\alpha$ 与 $\wt{\sigma}$。这三个提升连同 $\alpha$ 的一条半无穷射线提升 $\wt{\alpha}$，合成一个四边形 $\wt{\alpha},\wt{\delta}_\alpha,\wt{\beta}_\alpha,\wt{\sigma}$，其内角为 $\pi/2,\pi/2,\pi/2$，以及 $\wt{\sigma}$ 与 $\wt{\alpha}$ 之间的角 $0$。由四边形的双曲三角学 [Bus92] 可得
\[
\tanh\ell_{\beta_\alpha}=\operatorname{sech}\delta_\alpha=\frac{\sinh(\ell_\beta/2)\sinh(\ell_\alpha/2)}{\cosh(\ell_\lambda/2)+\cosh(\ell_\beta/2)\cosh(\ell_\alpha/2)}.
\]
补区间的长度为 $2\ell_{\beta_\alpha}=\ell_\beta-\mathcal{R}(\ell_\beta,\ell_\lambda,\ell_\alpha)$。主间隙长度的公式现在由一般关系 $\mathcal{R}(x,y,z)+\mathcal{R}(x,z,y)=x+\mathcal{D}(x,y,z)$ 推出。对双边界袖口裤子，主间隙长度要加到补区间长度上。结果为 $\mathcal{R}(\ell_\beta,\ell_\alpha,\ell_\lambda)$。
\end{proof}

% [Tutor Note: p.19 的证明里有一处可疑记号（已用 900 dpi 裁图核对印刷）：原文是"线段 \(\beta_\alpha\) 的一端
%  有测地线 \(\delta_\beta\) 出发"，但同一句紧接着的四边形四条边写作 \(\wt{\alpha},\wt{\delta}_\alpha,\wt{\beta}_\alpha,\wt{\sigma}\)，
%  而 \(\beta_\alpha\) 正是被 \(\delta_\alpha\) 平分出来的半区间，故该端点出发的应是 \(\delta_\alpha\)；
%  再者这条裤子只有 \(\beta,\alpha,\lambda\) 三条边界，不存在"\(\delta_\beta\)"这样的对象，
%  且 \(\delta_\beta\) 在全篇再未出现。判断依据：①四边形要闭合，四边必须首尾相接，缺了 \(\delta_\alpha\) 就拼不上；
%  ②凸性/切角条件 \(\pi/2,\pi/2,\pi/2,0\) 只与 \(\delta_\alpha,\beta_\alpha,\sigma,\alpha\) 这四条边相容。
%  按"正文照原样译、批注指正"的规则，正文保留 \(\delta_\beta\)；此处标出：极可能是 \(\delta_\alpha\) 之误。]

% [Tutor Note: p.19 的 \(\tanh\ell_{\beta_\alpha}=\operatorname{sech}\delta_\alpha\) 可用裤子三角学独立复核，
%  这也是验证 (2) 是否抄对的最好办法。设 \(\delta_\alpha\) 为连接 \(\beta\) 与 \(\alpha\) 的正交测地线的长度，则
%  \(\cosh\delta_\alpha=\dfrac{\cosh(\ell_\beta/2)\cosh(\ell_\alpha/2)+\cosh(\ell_\lambda/2)}{\sinh(\ell_\beta/2)\sinh(\ell_\alpha/2)}\)，
%  取倒数即得书中的那个分式（分母为 \(\cosh(\ell_\lambda/2)+\cosh(\ell_\beta/2)\cosh(\ell_\alpha/2)\)）。
%  另一条路线：把 \(\mathcal R\) 的定义代入 \(2\ell_{\beta_\alpha}=\ell_\beta-\mathcal R(\ell_\beta,\ell_\lambda,\ell_\alpha)\)
%  得 \(e^{2\ell_{\beta_\alpha}}=\big(\cosh(\ell_\lambda/2)+\cosh((\ell_\beta+\ell_\alpha)/2)\big)\big/\big(\cosh(\ell_\lambda/2)+\cosh((\ell_\beta-\ell_\alpha)/2)\big)\)，
%  再用 \(\tanh=\frac{N-D}{N+D}\) 与和差化积 \(\cosh a-\cosh b=2\sinh\frac{a+b}{2}\sinh\frac{a-b}{2}\)，
%  得到的正是同一分式——两条路线互相印证，同时验证了 \(\mathcal R\) 的三个变量顺序。
%  记号上注意：原文把"从 \(\beta\) 到 \(\alpha\) 的正交边界测地线"与其长度都写作 \(\delta_\alpha\)（\(\operatorname{sech}\delta_\alpha\) 里指的是长度），
%  这不是笔误，而是作者一贯的长度/测地线同名混用。
%  数值自检（可直接验收）：取 \((x,y,z)=(1,2,3)\)，由 (2) 算得 \(\mathcal R(1,2,3)\approx0.45823\)、\(\mathcal R(1,3,2)\approx0.69846\)，
%  和为 \(1.15668\)；再由 (1)(2) 算得 \(\mathcal D(1,2,3)\approx0.15668\)，故 \(x+\mathcal D=1.15668\)，
%  两者到小数点后五位完全相同，验证了证明中用到的一般关系 \(\mathcal R(x,y,z)+\mathcal R(x,z,y)=x+\mathcal D(x,y,z)\)。]

\para{第 3 讲：余体积公式与递推。}

主要步骤是应用长度恒等式，把映射类群的作用约化为更小映射类群的作用，从而把体积表达为某个长度水平集上的积分。结果就是低维体积函数之积的积分——即递推。

该方法通过计算亏格 $1$、一条边界的体积来说明。长度恒等式为
\[
L=\sum_{\alpha\ \text{简单}}\mathcal{D}(L,\ell_\alpha,\ell_\alpha).
\]
引入 $\Stab(\alpha)\subset\MCG$，即 $\MCG$ 作用在自由同伦类上的稳定化子。环面是一条椭圆曲线，其万有覆盖为 $\C$，带对合 $z\to-z$，它稳定覆盖变换格。该对合作用在环面上，也作用在带一个穿刺点或一条边界的环面上。对每条简单闭测地线的自由同伦类，该对合反转定向；而稳定化子 $\Stab(\alpha)$ 是 Dehn 扭转与对合生成的 $\Z/2\Z$ 子群的半直积。对合在 Teichmüller 空间上作用平凡。（环面是 $\MCG$ 在 $\cT$ 上的作用不忠实的例外情形之一。下面我们也会讨论环面情形，那里涉及一个重数。）一次 Dehn 扭转在 FN 坐标下作用在 Teichmüller 空间上为 $(\ell,\tau)\to(\ell,\tau+\ell)$。扇区 $\{0\le\tau<\ell\}$ 是 $\Stab(\alpha)$ 作用的基本域。映射类 $h\in\MCG$ 作用在测地长度函数上为 $\ell_\alpha\circ h^{-1}=\ell_{h(\alpha)}$。

% [Tutor Note: 末句 \(\ell_\alpha\circ h^{-1}=\ell_{h(\alpha)}\) 是全篇反复使用的"等变"约定，易与
%  \(\ell_\alpha\circ h=\ell_{h^{-1}(\alpha)}\) 混同，核对办法是取 \(h\) 为 Dehn 扭转：两条式子给出同一结论，
%  但只有 \(h^{-1}\) 在右的形式才能与 p.19 下行 \((\ell,\tau)\to(\ell,\tau+\ell)\) 的"左作用"约定相容
%  （模空间是 \(\cT/\MCG\)，故群作用从左边复合 \(h\cdot x\)，函数被拉回时出现 \(h^{-1}\)）。
%  另可核验扇区 \(\{0\le\tau<\ell\}\)：一次 Dehn 扭转把 \(\tau\) 加 \(\ell\)，故该半开区间的长度恰为 \(\ell\)，
%  其面积 \(\int_0^\ell\dd\tau=\ell\) 与 \(2\ell\) 边界的双曲面积公式一致，是取基本域的一个即时自检。]
